\lyseis{Β1.9}{Θέσεις ευθειών στο επίπεδο}
\sol{1} α) παράλληλες \quad β) τεμνόμενες \quad γ) παράλληλες \quad δ) τεμνόμενες (αν είναι διαφορετικές).
\sol{2} Ελεύθερη κατασκευή με γνώμονα και χάρακα.
\sol{3} Από σημείο εκτός ευθείας: μία. Από σημείο πάνω της: καμία παράλληλη (η ευθεία ταυτίζεται με τον εαυτό της).
\sol{4} $0$ (παράλληλες), $1$ (από το ίδιο σημείο), $2$ (δύο παράλληλες και μία που τις τέμνει), $3$ (ανά δύο διαφορετικά σημεία).
\sol{5} Για $4$ ευθείες το πολύ $\frac{4\cdot3}{2}=6$ σημεία. Για $\nu$ ευθείες το πολύ $\frac{\nu(\nu-1)}{2}$.
\sol{6} Είναι παράλληλες: $\varepsilon_1\parallel\varepsilon_3$.
\sol{7} α) Σωστό (στο ίδιο επίπεδο). β) Λάθος. γ) Σωστό. δ) Λάθος: περνά μία.
\sol{8} Βάζουμε παλούκια στο ίδιο μήκος με τη σειρά, μετράμε $60$ cm κάθετα από τη μία σειρά και τεντώνουμε το σχοινί ακριβώς παράλληλα.
\sol{9} $2$ ζεύγη παράλληλων πλευρών και $4$ ζεύγη κάθετων.
\sol{10} Το πολύ $\frac{5\cdot4}{2}=10$ σημεία.
\sol{11} Αν οι ράγες συναντιούνταν, το τρένο δεν θα χωρούσε. Η απόσταση πρέπει να μένει σταθερή.
\sol{12} Έχεις $3$ σημεία τομής.
\sol{13} Από δύο σημεία περνά μία ευθεία (δεν χρειάζεται άλλη προϋπόθεση). Η παράλληλη χρειάζεται και σημείο και ευθεία: είναι δηλαδή ανεξάρτητο, ισχύει ως ξεχωριστό αξίωμα.
\sol{14} Ναι, αν οι ευθείες δεν είναι παράλληλες συναντιούνται κάπου πέρα από το χαρτί.
\sol{15} Είναι παράλληλες.
\sol{16} Όχι: το αξίωμα του Ευκλείδη λέει ότι υπάρχει μία. Αν υπήρχαν δύο, η γεωμετρία θα ήταν διαφορετική.
\sol{17} Ελεύθερη απάντηση.
\sol{18} Ελεύθερη απάντηση.
